\chapter{مجموعه‌آیتم‌های پرتکرار و قوانین انجمنی}
\label{ch:frequent-itemsets}

\section{مقدمه و مدل سبد بازار \texorpdfstring{\enparen{Market-Basket Model}}{(Market-Basket Model)}}
یکی از کلاسیک‌ترین و در عین حال پرکاربردترین الگوها در کلان‌داده، کشف وابستگی‌ها میان اقلامی است که معمولاً به صورت هم‌زمان رخ می‌دهند. این چارچوب تحلیلی تحت عنوان \textbf{مدل سبد بازار}\footnote{\lr{Market-Basket Model}} شناخته می‌شود.

در مدل سبد بازار:
\begin{itemize}
    \item مجموعه‌ای بزرگ از \textbf{آیتم‌ها}\footnote{\lr{Items}} وجود دارد (مانند کالاها در فروشگاه آمازون).
    \item مجموعه‌ای بسیار بزرگ از \textbf{سبدها}\footnote{\lr{Baskets}} وجود دارد که هر سبد، زیرمجموعه‌ای کوچک از آیتم‌ها است (مانند سبد خرید یک مشتری در یک بار مراجعه).
    \item فرض بر این است که تعداد کل آیتم‌ها ممکن است در حد میلیون‌ها، و تعداد سبدها در حد صدها میلیون باشد، اما تعداد آیتم‌های درون هر سبد معمولاً اندک (مثلاً کمتر از ۲۰ قلم) است.
\end{itemize}

کاربردهای این مدل فراتر از خرده‌فروشی است:
\begin{itemize}
    \item \textbf{کاوش مفاهیم و کلمات هم‌نشین}: آیتم‌ها کلمات و سبدها اسناد متنی هستند (کشف عبارات مرکب).
    \item \textbf{زیست‌انفورماتیک}: آیتم‌ها نشانگرهای زیستی و سبدها پرونده ژنتیکی بیماران مبتلا به یک بیماری خاص هستند.
    \item \textbf{تحلیل سرقت ادبی}: آیتم‌ها جملات و سبدها مقالات هستند.
\end{itemize}

\section{تعاریف پایه: پشتیبانی، اطمینان و قوانین انجمنی}

\begin{definition}[پشتیبانی یک مجموعه‌آیتم]
\textbf{پشتیبانی}\footnote{\lr{Support}} یک مجموعه‌آیتم $I$ که با $s(I)$ نشان داده می‌شود، برابر است با تعداد سبدهایی که $I$ زیرمجموعه‌ای از آن‌ها است:
\begin{equation}
s(I) = \Big| \{ B \in \text{Baskets} : I \subseteq B \} \Big|
\end{equation}
اگر یک \textbf{آستانه پشتیبانی}\footnote{\lr{Support Threshold}} مانند $s$ داشته باشیم، هر مجموعه‌آیتم $I$ که شرط $s(I) \ge s$ را ارضا کند، یک \textbf{مجموعه‌آیتم پرتکرار}\footnote{\lr{Frequent Itemset}} نامیده می‌شود.
\end{definition}

\begin{definition}[قوانین انجمنی، اطمینان و بهره]
یک \textbf{قانون انجمنی}\footnote{\lr{Association Rule}} به فرم $I \to j$ بیان می‌شود؛ یعنی «اگر سبدی شامل مجموعه‌آیتم $I$ باشد، آنگاه این سبد با احتمال بالایی شامل آیتم $j$ نیز خواهد بود».
\begin{itemize}
    \item \textbf{اطمینان \enparen{Confidence}}: نسبت احتمال وقوع توأم به احتمال وقوع شرط:
    \begin{equation}
    \text{conf}(I \to j) = \frac{s(I \cup \{j\})}{s(I)}
    \end{equation}
    
    \item \textbf{بهره \enparen{Interest}}: تفاوت میان اطمینان قانون و احتمال وقوع تصادفی آیتم $j$ در کل سبدها:
    \begin{equation}
    \text{interest}(I \to j) = \text{conf}(I \to j) - P(j) = \text{conf}(I \to j) - \frac{s(\{j\})}{|Baskets|}
    \end{equation}
    قوانینی جذاب و ارزشمند هستند که بهره آن‌ها مقدار مثبت بزرگی داشته باشد (همبستگی مثبت قوی).
\end{itemize}
\end{definition}

\begin{example}[محاسبه اطمینان و بهره]
فرض کنید در یک فروشگاه آنلاین با $10{,}000$ سبد:
\begin{itemize}
    \item سبدهای حاوی پوشک: $s(\text{پوشک}) = 1000$ $(P(\text{پوشک}) = 0.1)$.
    \item سبدهای حاوی شیر خشک: $s(\text{شیر}) = 2000$ $(P(\text{شیر}) = 0.2)$.
    \item سبدهای حاوی هر دو: $s(\{\text{پوشک}, \text{شیر}\}) = 800$.
\end{itemize}
برای قانون $\text{پوشک} \to \text{شیر}$:
\begin{align}
\text{conf}(\text{پوشک} \to \text{شیر}) &= \frac{800}{1000} = 0.8 \quad (80\%) \\
\text{interest}(\text{پوشک} \to \text{شیر}) &= 0.8 - 0.2 = +0.6
\end{align}
چون بهره مثبت و بالاست، این قانون همبستگی قوی خرید همزمان را اثبات می‌کند.
\end{example}

\section{اصل یکنوایی و الگوریتم A-Priori}
چالش محاسباتی اصلی در یافتن مجموعه‌آیتم‌های پرتکرار، رشد نمایی ترکیبات است. اگر $N$ آیتم داشته باشیم، تعداد کل زیرمجموعه‌های ممکن $2^N$ است.
کلید غلبه بر این انفجار ترکیبیاتی، \textbf{اصل یکنوایی}\footnote{Monotonicity Property یا Apriori Principle} است.

\begin{theorem}[اصل یکنوایی پشتیبانی]
اگر یک مجموعه‌آیتم $I$ پرتکرار باشد، آنگاه تمامی زیرمجموعه‌های آن نیز قطعاً پرتکرار هستند.
به بیان عکس (عکس نقیض): \textbf{اگر مجموعه‌آیتمی نامتکرار باشد، هیچ اَبَرمجموعه‌ای از آن نمی‌تواند پرتکرار باشد}.
\begin{equation}
J \subseteq I \implies s(I) \le s(J)
\end{equation}
\end{theorem}

\subsection{مراحل اجرای الگوریتم A-Priori}
الگوریتم A-Priori در گذرها\footnote{\lr{Passes}} متوالی بر روی داده‌های سبدها عمل می‌کند:

\begin{algorithmbox}[الگوریتم استاندارد A-Priori~\cite{agrawal1994fast}]
\begin{itemize}
    \item \textbf{گذر اول \enparen{Pass 1}}:
    \begin{enumerate}
        \item با یک بار خواندن سبدها از روی دیسک، فراوانی هر آیتم منفرد شمرده می‌شود.
        \item آیتم‌هایی که فراوانی آن‌ها $\ge s$ باشد، مجموعه آیتم‌های ۱-تایی پرتکرار $(L_1)$ را تشکیل می‌دهند.
    \end{enumerate}
    
    \item \textbf{تولید نامزدهای ۲-تایی $(C_2)$}:
    تنها جفت‌هایی $\{i, j\}$ به عنوان نامزد در نظر گرفته می‌شوند که هر دوی $i \in L_1$ و $j \in L_1$ باشند.
    
    \item \textbf{گذر دوم \enparen{Pass 2}}:
    سبدها مجدداً از دیسک خوانده می‌شوند. برای هر سبد، تمامی جفت‌های درون آن که در $C_2$ عضویت دارند پیدا شده و شمارنده آن‌ها یکی افزایش می‌یابد. در پایان گذر دوم، جفت‌های با شمارش $\ge s$ به عنوان $L_2$ انتخاب می‌شوند.
    
    \item \textbf{گذر $k$-ام (Pass $k$)}:
    مجموعه نامزدهای $C_k$ از روی $L_{k-1}$ تولید می‌شود، سبدها خوانده شده و $L_k$ به دست می‌آید. این فرایند تا زمانی ادامه می‌یابد که هیچ نامزد جدیدی تولید نشود.
\end{itemize}
\end{algorithmbox}

\section{الگوریتم‌های بهینه‌سازی مصرف حافظه}
در الگوریتم A-Priori، گلوگاه اصلی حافظه در **گذر دوم** (شمارش تعداد جفت‌های نامزد $C_2$) رخ می‌دهد. الگوریتم‌های زیر برای حل این گلوگاه ابداع شدند:

\subsection{الگوریتم پارک-چن-یو \texorpdfstring{\enparen{PCY}}{(PCY)}}
\textbf{مشاهده کلیدی}: در گذر اول الگوریتم A-Priori، تنها کاری که انجام می‌شود شمارش آیتم‌های تکی است که حافظه بسیار ناچیزی نیاز دارد؛ در نتیجه بخش اعظم حافظه اصلی کامپیوتر در گذر اول بیکار می‌ماند!

\textbf{ایده الگوریتم PCY}:
\begin{enumerate}
    \item در گذر اول، از حافظه آزاد باقیمانده به عنوان یک \textbf{جدول درهم‌سازی بزرگ از سطل‌ها}\footnote{\lr{Hash Table of Buckets}} استفاده می‌کنیم.
    \item هنگام خواندن هر سبد در گذر اول، علاوه بر شمارش آیتم‌های تکی، به ازای هر جفت $\{i, j\}$ در سبد، مقدار درهم‌سازی $h(i, j)$ را محاسبه کرده و شمارنده آن سطل را یک واحد افزایش می‌دهیم.
    \item در پایان گذر اول، اگر مقدار یک سطل کمتر از آستانه $s$ باشد، \textbf{امکان ندارد هیچ جفت پرتکراری به آن سطل نگاشت شده باشد}.
    \item جدول درهم‌سازی به یک \textbf{نقشه بیتی}\footnote{\lr{Bitmap}} فشرده تبدیل می‌شود (۱ بیت برای هر سطل: ۱ اگر مجموع سطل $\ge s$ باشد، وگرنه ۰).
    \item \textbf{در گذر دوم}: یک جفت $\{i, j\}$ به عنوان نامزد شمارش می‌شود اگر و تنها اگر:
    \begin{equation}
    i \in L_1 \quad \text{و} \quad j \in L_1 \quad \text{و} \quad \text{Bitmap}[h(i, j)] == 1
    \end{equation}
\end{enumerate}
این شرط سه‌گانه، تعداد جفت‌های نامزد در گذر دوم را به شدت کاهش می‌دهد.

\subsection{الگوریتم‌های چندمرحله‌ای و چنددرهم‌سازی \texorpdfstring{\enparen{Multistage \& Multihash}}{(Multistage \& Multihash)}}
\begin{itemize}
    \item \textbf{الگوریتم چندمرحله‌ای \enparen{Multistage}}: یک گذر میانی (گذر ۱.۵) اضافه می‌کند که در آن، تنها جفت‌هایی که از فیلتر سطل اول عبور کرده‌اند با یک تابع درهم‌سازی مستقل دوم هش می‌شوند؛ بدین ترتیب نقشه بیتی دوم تولید شده و تعداد نامزدهای گذر نهایی باز هم کاهش می‌یابد.
    \item \textbf{الگوریتم چنددرهم‌سازی \enparen{Multihash}}: در همان گذر اول، حافظه را به دو یا چند جدول درهم‌سازی مستقل با دو تابع هش مجزا تقسیم می‌کند.
\end{itemize}

\section{الگوریتم‌های با تعداد گذر محدود (یک یا دو گذره)}
در کلان‌داده، خواندن مکرر ترابایت‌ها داده از دیسک زمان‌بر است. الگوریتم‌های زیر یافتن مجموعه‌آیتم‌های پرتکرار را در ۱ یا ۲ گذر میسر می‌سازند:

\subsection{الگوریتم ساواسره-امینسکی-ناواته \texorpdfstring{\enparen{SON}}{(SON)}}
الگوریتم SON ساختاری کاملاً موازی دارد و بهترین الگوریتم برای پیاده‌سازی در \textbf{MapReduce} است:

\begin{algorithmbox}[الگوریتم SON در MapReduce]
\begin{enumerate}
    \item \textbf{گذر اول (گام MapReduce اول)}:
    \begin{itemize}
        \item داده‌ها به $m$ قطعه مساوی \enparen{Chunk} تقسیم می‌شوند.
        \item هر مپر یک قطعه را خوانده و با استفاده از الگوریتم A-Priori یا PCY، مجموعه‌آیتم‌های پرتکرار \textbf{محلی} با آستانه متناسب $s/m$ را پیدا می‌کند.
        \item \textbf{لم بنیادین}: هر مجموعه‌آیتمی که در کل داده‌ها پرتکرار باشد، \textbf{حداقل در یکی از $m$ قطعه} باید دارای پشتیبانی $\ge s/m$ باشد.
        \item کاهنده اجتماع تمامی مجموعه‌آیتم‌های پرتکرار محلی را به عنوان نامزدهای جهانی منتشر می‌سازد.
    \end{itemize}
    
    \item \textbf{گذر دوم (گام MapReduce دوم)}:
    مپرها سبدها را خوانده و پشتیبانی دقیق نامزدهای جهانی را شمارش می‌کنند؛ کاهنده نامزدهایی را که مجموعاً $\ge s$ هستند به عنوان پاسخ قطعی و نهایی برمی‌گزیند.
\end{enumerate}
\end{algorithmbox}
ویژگی منحصربه‌فرد SON: \textbf{دقیقاً در ۲ گذر اجرا می‌شود و هیچ منفی کاذب یا مثبت کاذبی تولید نمی‌کند!}

\subsection{الگوریتم تویونن \texorpdfstring{\enparen{Toivonen's Algorithm}}{(Toivonen's Algorithm)} و مرز منفی}
الگوریتم Toivonen با یک گذر بر روی نمونه و یک گذر بر روی کل داده‌ها عمل می‌کند:
\begin{enumerate}
    \item نمونه‌ای تصادفی با اندازه کسری $p$ از سبدها برداشته می‌شود.
    \item مجموعه‌آیتم‌های پرتکرار در نمونه با آستانه‌ای اندکی کمتر از $p \cdot s$ استخراج می‌گردند.
    \item \textbf{مرز منفی \enparen{Negative Border}} تشکیل می‌شود: مجموعه‌آیتم‌هایی که خود در نمونه پرتکرار نشده‌اند، اما \textbf{تمامی زیرمجموعه‌های بلافصل آن‌ها پرتکرار بوده‌اند}.
    \item در گذر دوم بر روی کل داده‌های اصلی، پشتیبانی تمامی آیتم‌های پرتکرار نمونه و تمامی آیتم‌های مرز منفی به طور هم‌زمان شمرده می‌شود.
    \item \textbf{شرط موفقیت}: اگر هیچ‌یک از اعضای مرز منفی در داده‌های اصلی به آستانه $s$ نرسند، نتایج نمونه برای داده‌های اصلی \textbf{کاملاً دقیق و بدون خطا} است! (اگر عضوی از مرز منفی پرتکرار شود، آزمون باطل شده و با نمونه جدید تکرار می‌گردد).
\end{enumerate}

\begin{examnote}
نکات امتحانی پرکاربرد:
\begin{enumerate}
    \item \textbf{فرمول آدرس‌دهی ماتریس مثلثی}: برای ذخیره شمارنده‌های جفت‌های نامزد در یک آرایه یک‌بعدی فشرده (بدون اتلاف فضای قطر و بخش بالایی)، درایه $(i, j)$ با شرط $i < j$ در خانه زیر ذخیره می‌شود:
    \begin{equation}
    k = (i - 1) \left( n - \frac{i}{2} \right) + j - i
    \end{equation}
    مقایسه روش ماتریس مثلثی در برابر سه‌تایی‌ها \enparen{Triples}: هر خانه ماتریس مثلثی ۴ بایت (یک عدد صحیح) فضا می‌گیرد؛ اما روش سه‌تایی‌ها برای هر جفت نیازمند ۱۲ بایت $((i, j, \text{count}))$ است. بنابراین اگر بیش از یک‌سوم کل جفت‌های ممکن نامزد باشند، ماتریس مثلثی فشرده‌تر است؛ اما اگر جفت‌های نامزد بسیار تنک باشند، روش سه‌تایی‌ها کاراتر خواهد بود.
    \item \textbf{اثبات صفر بودن منفی‌های کاذب در الگوریتم SON}:
    اگر مجموعه‌آیتم $I$ در کل داده‌ها دارای پشتیبانی حداقل $s$ باشد، غیرممکن است که در تمام $m$ قطعه دارای پشتیبانی کمتر از $s/m$ باشد؛ زیرا در این صورت مجموع کل پشتیبانی اکیداً کمتر از $m \times (s/m) = s$ می‌شد که تناقض است. بنابراین $I$ قطعاً حداقل در یک قطعه کشف شده و هیچ آیتم پرتکراری از دست نمی‌رود.
    \item \textbf{نقش حیاتی مرز منفی در الگوریتم Toivonen}: اگر کوچک‌ترین آیتمی که در کل داده‌ها به حد نصاب رسیده اما در نمونه کشف نشده وجود داشته باشد، طبق اصل یکنوایی تمام زیرمجموعه‌های بلافصل آن باید در نمونه کشف شده باشند؛ پس آن آیتم حتماً عضوی از مرز منفی است و در گذر دوم شمرده می‌شود.
\end{enumerate}
\end{examnote}

\begin{summarybox}[جمع‌بندی فصل ششم]
\begin{itemize}
    \item مدل سبد بازار همبستگی‌های هم‌نشینی را از طریق معیارهای پشتیبانی، اطمینان و بهره ارزیابی می‌کند.
    \item اصل یکنوایی با حذف اَبَرمجموعه‌های اقلام نامتکرار، مانع انفجار ترکیبیاتی می‌شود.
    \item الگوریتم PCY با به کارگیری حافظه بلااستفاده برای سطل‌های درهم‌سازی، گلوگاه گذر دوم را مرتفع می‌سازد.
    \item الگوریتم SON چارچوبی بهینه برای موازی‌سازی کاوش در بستر MapReduce تنها در دو گذر است.
    \item الگوریتم Toivonen به کمک مرز منفی، بررسی صحت نتایج نمونه‌گیری را تضمین می‌نماید.
\end{itemize}
\end{summarybox}
